% Part: history
% Chapter: biographies 
% Section: kurt-goedel

\documentclass[../../../include/open-logic-section]{subfiles}

\begin{document}

\olfileid{his}{bio}{god}

\olsection{కుర్ట్ గ్యోడెల్}

\olphoto{goedel-kurt}{కుర్ట్ గ్యోడెల్}

కుర్ట్ గ్యోడెల్ (ఉచ్చారణ: గెర్-డల్) 1906 ఏప్రిల్ 28న ఆస్ట్రో-హంగేరియన్ సామ్రాజ్యంలోని బ్రూన్‌లో (ప్రస్తుతం చెక్ రిపబ్లిక్‌లోని బ్ర్నో) జన్మించారు. చిన్న కుర్టెలె ప్రశ్నించే స్వభావం, తెలివితేటల వల్ల కుటుంబ సభ్యులు ఆయనను తరచుగా ``Der kleine Herr Warum'' (చిన్న ``ఎందుకు'' గారు) అని పిలిచేవారు. ప్రాథమిక పాఠశాల నుంచే ఆయన చదువులో రాణించారు; గణితంలో మాత్రమే అత్యున్నత శ్రేణి కంటే తక్కువ మార్కు పొందారు. అనారోగ్యం వల్ల గ్యోడెల్ తరచూ పాఠశాలకు హాజరుకాలేకపోయేవారు; శారీరక విద్య నుంచి మినహాయింపు పొందారు. చిన్నతనంలో ఆయనకు రుమాటిక్ జ్వరం ఉన్నట్లు నిర్ధారణ అయింది. వైద్య పరీక్షలు విరుద్ధంగా చెప్పినా, అది తన గుండెను శాశ్వతంగా ప్రభావితం చేసిందని జీవితాంతం నమ్మారు.

గ్యోడెల్ 1924లో వియన్నా విశ్వవిద్యాలయంలో చదువు మొదలుపెట్టి,~1929లో డాక్టరేట్ పూర్తిచేశారు. మొదట భౌతికశాస్త్రం చదవాలనుకున్నారు; కానీ తత్వవేత్త రూడోల్ఫ్ కార్నాప్ ప్రభావం కొంత కారణంగా, ఆయన ఆసక్తి త్వరలో గణితం, ముఖ్యంగా తర్కశాస్త్రం వైపు మళ్లింది. హాన్స్ హాన్ పర్యవేక్షణలో రాసిన ఆయన పరిశోధన గ్రంథం, తాదాత్మ్యంతో కూడిన ప్రథమ శ్రేణి విధేయ తర్కం సంపూర్ణత సిద్ధాంతాన్ని నిరూపించింది \citep{Godel1929}. కేవలం ఏడాది తరువాత ఆయనకు అత్యంత పేరు తెచ్చిన ఫలితాలు---మొదటి, రెండవ అసంపూర్ణత సిద్ధాంతాలు---సాధించారు (\citealt{Godel1931}లో ప్రచురితం). వియన్నాలో ఉన్నప్పుడు గ్యోడెల్, శాస్త్రీయ దృక్కోణమున్న తత్వవేత్తల సమూహమైన వియన్నా వలయంలో చురుకుగా పాల్గొన్నారు. అందులో కార్నాప్ కూడా ఉన్నారు; ఆయన కృషిపై గ్యోడెల్ ఫలితాల ప్రభావం ప్రత్యేకంగా పడింది.

1938లో గ్యోడెల్ అడెలె నింబుర్స్కీని వివాహం చేసుకున్నారు. ఆయన తల్లిదండ్రులకు ఇది నచ్చలేదు: ఆమె ఆయన కంటే ఆరు సంవత్సరాలు పెద్దదే కాక అప్పటికే విడాకులు తీసుకున్నారు; రాత్రి వినోదశాలలో నర్తకిగా కూడా పనిచేసేవారు. అయితే సామాజిక ఒత్తిళ్లు గ్యోడెల్‌పై ప్రభావం చూపలేదు; ఆయన మరణించే వరకు వారు సంతోషంగా దాంపత్య జీవితం గడిపారు.

1938లో నాజీ జర్మనీ ఆస్ట్రియాను విలీనం చేసుకున్న తరువాత గ్యోడెల్, అడెలె అమెరికా సంయుక్త రాష్ట్రాలకు వలస వెళ్లారు. అక్కడ న్యూజెర్సీలోని ప్రిన్స్‌టన్‌లో ఉన్న ఇన్‌స్టిట్యూట్ ఫర్ అడ్వాన్స్‌డ్ స్టడీలో ఆయన పదవి చేపట్టారు. అంతర్ముఖ స్వభావం, విచిత్ర ప్రవర్తన ఉన్నప్పటికీ, ప్రిన్స్‌టన్‌లో ఆయన గడిపిన కాలం సహకారపూర్వకంగా, ఫలప్రదంగా సాగింది. సమితి సిద్ధాంతం, తత్వశాస్త్రం, భౌతికశాస్త్రంలో వ్యాసాలు ప్రచురించారు. ముఖ్యంగా, అదే సంస్థలోని సహోద్యోగి ఆల్బర్ట్ ఐన్‌స్టీన్‌తో ఆయనకు గాఢ స్నేహం ఏర్పడింది.

చివరి సంవత్సరాలలో గ్యోడెల్ మానసిక ఆరోగ్యం క్షీణించింది. 1977లో ఆయన భార్య ఆసుపత్రిలో చేరడంతో ఆయనకు భోజనం వండిపెట్టలేకపోయారు. జీవితమంతా మానసిక ఆరోగ్య సమస్యలతో బాధపడిన ఆయన, తరువాత తీవ్రమైన అపనమ్మకానికి లోనయ్యారు. ఆహారంలో విషం కలుపుతారనే ప్రాణభయంతో గ్యోడెల్ తినడానికి నిరాకరించారు. 1978 జనవరి 14న ప్రిన్స్‌టన్‌లో ఆకలితో మరణించారు.

\begin{reading}
గ్యోడెల్ సమగ్ర జీవిత చరిత్రకు \citet{Dawson1997} చూడండి. ఇతర జీవిత చరిత్ర రచనలకు, తర్కశాస్త్రం, తత్వశాస్త్రాలకు గ్యోడెల్ చేసిన కృషిపై వ్యాసాలకు \citet{Wang1990}, \citet{Baaz2011}, \citet{Takeuti2003}, \citet{Sigmund2007} చూడండి.

గ్యోడెల్ పీహెచ్‌డీ పరిశోధన గ్రంథం మూల జర్మన్‌లో అందుబాటులో ఉంది \citep{Godel1929}. అసంపూర్ణత సిద్ధాంతాల మూల పాఠ్యం \citep{Godel1931}. గ్యోడెల్ ప్రచురిత, అప్రచురిత రచనలన్నీ, ఎంపిక చేసిన ఉత్తరప్రత్యుత్తరాలు ఆయన ఆంగ్ల \emph{Collected Papers}లో లభిస్తాయి \citet{Godel1986,Godel1990}.

గ్యోడెల్ అసంపూర్ణత సిద్ధాంతాల విస్తృత అధ్యయనానికి \citet{Smith2013} చూడండి. ఆ సిద్ధాంతాలపై అనౌపచారిక తాత్విక చర్చకు మార్క్ లిన్సెన్‌మేయర్ పాడ్‌కాస్ట్ చూడండి \citep{Linsenmayer2014}.
\end{reading}

\end{document}
